\documentclass[11pt]{article}
\usepackage[utf8]{inputenc}
\usepackage{lmodern}
\usepackage[T1]{fontenc}
\usepackage[margin=1in]{geometry}
\usepackage{amsmath,amssymb}
\usepackage{booktabs}
\usepackage{caption}
\usepackage{microtype}
\usepackage[hidelinks]{hyperref}

\captionsetup{font=small,labelfont=bf}
\setlength{\parskip}{0.4em}
\setlength{\parindent}{0pt}
\newcommand{\rmean}{\langle r\rangle}
\newcommand{\chif}{\chi_{f}}

\title{\textbf{Multifractality as a sensing resource:\\ a $D_2$ figure of merit for Rosenzweig--Porter
systems,\\ and a lattice platform}}
\author{J. Arturo Ornelas Brand\\[2pt]
\small Independent Researcher, Aguascalientes, M\'exico\\
\small \texttt{arturoornelas62@gmail.com}}
\date{July 11, 2026 --- draft v0.1 (for review)}

\begin{document}
\maketitle

\begin{abstract}
\noindent Two established results have never been composed. The fidelity susceptibility of a quantum
state to a Hamiltonian parameter equals one quarter of its quantum Fisher information,
$F_Q = 4\chif$~\cite{liu2014,safranek2017}; and an exact calculation in the Rosenzweig--Porter (RP)
ensemble shows that eigenfunction fidelity susceptibility is \emph{enhanced} in the fractal
non-ergodic-extended phase~\cite{sak2022}. Composed, they make the eigenvector fractal dimension
$D_2$ --- until now a diagnostic of localization physics --- an operational \textbf{metrological
figure of merit}: along the RP fractal phase, at fixed distance from criticality and for field-like
(local) couplings to the sensed parameter, the finite-size enhancement exponent of the quantum
Fisher information over the ergodic reference is $1-D_2$. We confirm this in-house under
preregistered criteria on the generalized-RP ensemble with $D_2$ \emph{measured} from
inverse-participation-ratio scaling (not assumed): enhancement $=1-D_2$ within $0.15$ at $4/4$
interior points of the phase (deviations $0.009$--$0.078$; correlation $0.996$). Two further
preregistered experiments fix the criterion's boundaries: geometric fractality alone --- a
disorder-free Sierpinski-gasket lattice with certified eigenpairs and degenerate-multiplet
treatment --- shows \textbf{no effect} (enhancement $-0.12$, $|{\cdot}|<0.15$; the RP formula does
not extend), while critical power-law banded random matrices \textbf{extend the law} at all three
tested bandwidths (deviations $0.016$--$0.064$). The criterion is therefore: \emph{multifractal
eigenstates of random ensembles, sensed through local couplings} --- randomness, not fractal
geometry, is the missing ingredient, and an operator-structure scope law (diagonal inherits;
short-range inherits once its range fits the fractal support; collective and dense couplings are
$D_2$-blind) gives the design rule. On the platform side we construct --- because it cannot be
cited --- the minimal model that carries the physics onto cold atoms: an \emph{interacting
tilted$+$disordered lattice} in which the tilt is only the sensed parameter (the bare
Wannier--Stark ladder is integrable, so sourcing RP statistics from gravity would be immediately
refutable) and the dense RP-like spectrum comes from interactions and disorder. A preregistered
coexistence test on the interacting $t$--$V$ chain finds a textbook ergodic$\to$localized crossover
($\rmean$: $0.532\to0.386$) whose \emph{typical} susceptibility peaks at the crossover's ergodic
edge --- not mid-crossover --- with the predicted heavy $x^{-2}$ tail (Hill index $1.96$--$2.04$)
in the chaotic regime and estimator-dependent rare-event gains inside the crossover. The platform
statement is scoped accordingly, and one instrument caveat is stated up front: this attaches to
\emph{lattice/Bloch} gravimetry, not to free-fall light-pulse interferometry. Tier-2 throughout:
a shared random-matrix description with preregistered numerical evidence, not a hardware claim.
\end{abstract}

\section{Introduction}

Two sibling programs measure the same observable on different substrates. A physics program tests a
Rosenzweig--Porter (RP) hypothesis for freely vibrating plates --- proposing the hypothesis with a
preregistered eigenvector fractal dimension~\cite{ornelas2026hyp} and then executing its registered
numerical campaign on true operators~\cite{ornelas2026fem}. An AI program finds that a neural
network's fabricated Poisson$\to$GOE transition carries genuinely RP-multifractal Hessian
eigenvectors, $D_2 = 0.76$ within an architecture band $[0.56,0.80]$~\cite{ornelas2026rmt}. In both,
$D_2$ is the \emph{decisive diagnostic}: it separates a genuine multifractal phase from spectra that
merely imitate intermediate spacing statistics.

This note asks what $D_2$ is \emph{for}, beyond diagnosis. The metrology literature has recently
begun to name multifractal states ``a resource for quantum sensing''~\cite{hu2026}, and
criticality-enhanced sensing is an established framework~\cite{sahoo2024,liu2024me}; but no
standalone, quantitative figure of merit tied to the fractal dimension itself has been formulated
--- the notion is adjacent and nameable, yet the cell is empty. Meanwhile the two ingredients of
such a figure of merit have existed, separately, for years: the identity $F_Q = 4\chif$ between
quantum Fisher information and fidelity susceptibility~\cite{liu2014,safranek2017,hauke2016}, and
Kravtsov and collaborators' exact result that eigenfunction fidelity susceptibility is enhanced in
the RP fractal phase~\cite{sak2022} --- derived in the same ensemble the plate program
uses~\cite{kravtsov2015}.

We compose them, state the resulting figure of merit operationally, and then do what the parent
programs' discipline requires: preregister the readings, run the tests, and report the boundaries
of validity with the same weight as the confirmations. The result is a criterion --- multifractal
eigenstates of \emph{random} ensembles, sensed through \emph{local} couplings --- that survived one
designed attempt to break it (a disorder-free fractal lattice, where it correctly reports no
effect) and extended without adjustment to a third ensemble family. Part~1 of the contribution
constructs the minimal cold-atom platform on which the figure of merit is operational; Part~2 is
the figure of merit itself. We present Part~2 first, since it is load-bearing.

\section{The figure of merit}
\label{sec:fom}

\textbf{Definitions.} For a Hamiltonian $H(h) = H + h\,V$ and eigenstate $|n(h)\rangle$, the
fidelity susceptibility is
\[
\chif^{(n)} \;=\; \lim_{h\to0}\, \frac{-2\ln|\langle n(0)|n(h)\rangle|}{h^{2}}
\;=\; \sum_{m\neq n} \frac{|\langle m|V|n\rangle|^{2}}{(E_m-E_n)^{2}},
\]
and $F_Q = 4\chif$ is the quantum Fisher information of the eigenstate as a probe of
$h$~\cite{liu2014,safranek2017}: the Cram\'er--Rao-optimal sensitivity to the parameter, up to a
factor fixed by convention. Spectroscopy-accessible formulations of QFI are
standard~\cite{hauke2016}.

\textbf{The exact input.} In the RP ensemble --- $H = \mathrm{diag}(\epsilon_i) +
N^{-\gamma/2}W$ with GOE $W$ --- the fractal non-ergodic-extended phase at $1<\gamma<2$ has
eigenstates occupying $N^{D_2}$ of $N$ available sites, $D_2 = 2-\gamma$~\cite{kravtsov2015}.
Skvortsov, Amini and Kravtsov~\cite{sak2022} computed the eigenfunction fidelity susceptibility
exactly in this ensemble and found it \emph{enhanced} in the fractal phase: the mini-band structure
concentrates spectral weight at small denominators.

\textbf{The operational statement.} Let $\chif^{\rm typ}(N)$ be the typical (log-averaged)
mid-spectrum fidelity susceptibility at system size $N$ for a \emph{field-like} perturbation ---
diagonal in the localization basis, the analogue of a local field or a tilt --- and let
$\chif^{\rm erg}(N)$ be the same quantity for a dense (GOE-like) perturbation, which we show below
is $D_2$-blind and therefore serves as the in-situ ergodic reference. Define the enhancement
exponent $\alpha_{\rm enh}$ by $\chif^{\rm typ}/\chif^{\rm erg} \sim N^{\alpha_{\rm enh}}$. The
figure of merit is
\[
\boxed{\;\alpha_{\rm enh} \;=\; 1 - D_2\;}
\]
along the fractal phase, at fixed distance from its edges, with $D_2$ the \emph{measured}
inverse-participation-ratio scaling dimension of the same eigenstates. Smaller fractal dimension
--- more concentrated states --- buys polynomially more Fisher information per added system size,
relative to what an ergodic sensor of the same size would gain.

\section{Preregistered confirmation on measured $D_2$}
\label{sec:q1}

All experiments in this note follow the parent programs'
discipline~\cite{ornelas2026fem}: readings frozen in the experiment header before the run, verdicts
quoted as frozen, amendments dated. The registry and code are in the companion
repository.\footnote{\texttt{rp-qfi-sensing}; experiments \texttt{q01}--\texttt{q05}. Release
plans for the repository accompany the deposit of this note.}

\textbf{Design (Q1/Q1b).} Generalized-RP ensemble at interior points
$\gamma \in \{1.2, 1.4, 1.6, 1.8\}$ --- each at least $0.2$ from both phase edges, so criticality
is held at fixed distance --- with ergodic ($\gamma=0.5$) and localized ($\gamma=2.5$) controls;
$N = 256$--$4096$; $\chif$ per mid-spectrum eigenstate for a diagonal (field-like) perturbation and
for a normalized GOE perturbation; $D_2$ measured from IPR scaling, never assumed equal to
$2-\gamma$.

\textbf{Results.} Measured $D_2 = 0.854/0.608/0.397/0.168$ against the theoretical
$0.8/0.6/0.4/0.2$ (agreement $0.03$--$0.05$). The diagonal-coupling exponent obeys
$\alpha_{\rm diag} = 2-D_2$ within $0.03$--$0.06$, while the GOE-coupling exponent is flat at
$\alpha_{\rm GOE} \approx 1.00$ across the entire phase: \textbf{a dense perturbation is
$D_2$-blind}, which is what licenses it as the in-situ ergodic reference. The enhancement exponent
$\alpha_{\rm diag}-\alpha_{\rm GOE}$ equals $1-D_2$ within $0.06$ at all four interior points. A
second, fully frozen-gate run on an independent seed (Q1b; gate $|\alpha_{\rm enh}-(1-D_2)| \le
0.15$ at $\ge 3/4$ points, declared before the run) passes $4/4$ with deviations
$0.015/0.018/0.009/0.078$ and $\mathrm{corr}(\alpha_{\rm enh},\,1-D_2) = 0.996$; monotonicity
across the phase passes with total variation $0.709$. One scoping finding emerged and is promoted
to part of the claim: \emph{the figure of merit requires field-like coupling to the sensed
parameter} --- it is a statement about how multifractal support meets a local operator, not about
the spectrum alone.

\section{The boundaries: what multifractality is not enough}
\label{sec:scope}

\textbf{Geometric fractality alone fails (Q3b).} The sharpest referee question the composition
faces~\cite{sak2022,tokarczyk2026} is whether ``multifractality helps'' is generic. We built the
adversarial test: a disorder-free Sierpinski-gasket tight-binding lattice, whose eigenstates are
multifractal from geometry alone~\cite{sierpinski2025} --- no randomness, no RP mini-bands.
Instrumentation matters here and is reported because it changed a verdict: a first run was
invalidated by an eigensolver defect (LAPACK \texttt{dsyevd} via \texttt{numpy} returning
eigenvalues off by up to $31$ at $N=9843$) and by a degeneracy-exclusion design that starved the
statistics; the certified rerun (independent eigenvalue reference, \texttt{evr}-driver vectors,
agreement and residual gates per size, and degenerate-multiplet subspace rotation --- correct
degenerate perturbation theory) passes certification at all four sizes with full-window statistics.
Frozen verdict: enhancement $-0.123$, i.e.\ \textbf{no effect} ($|{\cdot}|<0.15$), and the RP
formula does not extend (deviation $0.816$ from $1-D_2$). The robust statement, given the gasket's
composition scatter, is the absence of any consistent enhancement.

\textbf{A third random family extends the law (Q5).} Critical power-law banded random
matrices~\cite{mirlin1996} at bandwidths $b = 0.5/1.5/4.0$: $\alpha_{\rm enh} = 1-D_2^{\rm meas}$
within $0.016$--$0.064$ at all three. Combined with Q1b (RP: yes) and Q3b (deterministic gasket:
no), the criterion is not ``RP specifically'' and not ``fractality generically'' but:
\[
\textbf{multifractal eigenstates of random ensembles, sensed through local couplings.}
\]
What the gasket lacked was randomness --- its exact deterministic multiplets carry no mini-band
level structure for a local operator to exploit.

\textbf{The suppression counterexample, reconciled.} In the 3D Anderson model, fidelity
susceptibility is \emph{sub}maximal at the multifractal critical point, and that suppression is
attributed to the fractal structure of critical states~\cite{tokarczyk2026}. This does not
contradict the figure of merit; it delimits it. The RP result lives on an extended fractal
\emph{phase} at fixed criticality distance; the Anderson result compares a single critical
\emph{point} against its own ergodic side. ``Multifractality helps'' is model-dependent --- which
is precisely why the criterion above names the random-ensemble phase, not multifractality at
large.

\section{The operator scope law}
\label{sec:q4}

Which couplings inherit the enhancement (Q4, frozen classification on gRP at
$\gamma \in \{1.4, 1.8\}$):

\begin{center}
\begin{tabular}{lll}
\toprule
coupling $V$ & verdict & reading\\
\midrule
diagonal (field-like) & \textbf{inherits} & at both $\gamma$; the figure of merit's home\\
banded, range 8 & \textbf{conditional} & $58\%$ at $\gamma=1.4$; full at $\gamma=1.8$\\
low-rank (collective) & blind & no $D_2$ dependence\\
dense (GOE-like) & blind & the ergodic reference of \S\ref{sec:q1}\\
\bottomrule
\end{tabular}
\end{center}

A short-range operator inherits the enhancement once its range fits inside the fractal support
scale; collective and dense couplings never do. The design rule for any platform is therefore:
\textbf{sense with local couplings}. (The tilt used in \S\ref{sec:platform} is field-like by
construction.)

\section{The platform: an interacting tilted$+$disordered lattice}
\label{sec:platform}

\textbf{Why constructed rather than cited.} A systematic, adversarially verified literature search
(documented in the program's research record) found each of the bridge's four required pieces
published separately --- gravity entering a lattice as a
Wannier--Stark tilt with $\partial E_n/\partial g$ a literal eigenvalue
derivative~\cite{zapata2001,wilkinson1996}; spectroscopic access to QFI~\cite{hauke2016};
$F_Q = 4\chif$ instantiated on tilted and quasiperiodic lattices~\cite{sahoo2024,sahoo2024b} ---
but never composed with RP-class level statistics, and with two load-bearing caveats confirmed
independently: the bare Wannier--Stark ladder is \emph{integrable} (an equidistant picket fence,
the opposite of a dense RP spectrum), and the existing cold-atom QFI enhancements are
criticality effects with no level-statistics content. The minimal defensible model is therefore an
\textbf{interacting, disordered, tilted lattice} (many-body Bose--Hubbard or interacting
Aubry--Andr\'e--Harper class), with a non-negotiable design rule: \emph{the tilt is only the sensed
parameter}; the RP-like statistics must come from interactions and disorder, never from gravity.

\textbf{The coexistence test (Q2).} The decisive preregistered question: does such a system show
its Poisson$\leftrightarrow$GOE crossover in the same window where the fidelity susceptibility to
the tilt is large? Instantiation: the interacting $t$--$V$ chain (spinless fermions, $J=V=1$) with
onsite disorder $W$ and the tilt as sensed parameter at $h=0$; $L=12$ (Hilbert dimension 924, 200
realizations) and $L=14$ (3432, 64); first realization of every cell certified against an
independent eigensolver. Results: a textbook crossover ($\rmean: 0.532 \to 0.386$ as $W: 0.5\to8$);
the \emph{typical} susceptibility peaks at $W^{*}=2$ where $\rmean = 0.520$ --- the crossover's
\textbf{ergodic edge}, not mid-crossover, so the frozen mid-crossover gate reads CHALLENGES and the
platform statement is re-scoped accordingly. The predicted GOE $x^{-2}$ heavy tail of the $\chif$
distribution is confirmed (Hill index $1.96$--$2.04$ across the chaotic regime); \emph{inside} the
crossover the tail index drops below $2$ ($1.73/1.54$): the mean susceptibility explodes on rare
resonances while the typical value collapses, so sensitivity there is estimator-dependent.

\textbf{The platform statement (as re-scoped, honestly).} Maximum \emph{typical} sensitivity to a
field-like parameter sits at the chaotic edge of the crossover; the crossover interior offers
heavy-tailed, rare-event gains that a mean-based protocol can harvest and a typical-based protocol
cannot. Whether the peak tracks the finite-size crossover point or pins to fixed disorder is
registered as the follow-up (Q2c, finer grid, three sizes).

\textbf{The instrument caveat, stated before it is asked.} A light-pulse (Kasevich--Chu) atom
interferometer reads gravity from a free-fall phase $\Delta\phi = k_{\rm eff}\,g\,T^{2}$, which is
not a spectral quantity. Everything above lives on \emph{lattice/Bloch gravimetry} --- the
Wannier--Stark/Bloch-oscillation class of instruments, a cousin of the free-fall interferometer,
not its own readout. The bridge from RP structure to a light-pulse phase, if one exists, is future
work and is not claimed here.

\section{Relation to the parent programs, and outlook}
\label{sec:outlook}

The figure of merit gives the parent programs' $D_2$ measurements an operational meaning.
The AI program's fabricated transition ($D_2 = 0.76$, band $[0.56,0.80]$~\cite{ornelas2026rmt})
and the designed plate assembly that lands in the same band at router-like coupling
($D_2 = 0.63$--$0.69$~\cite{ornelas2026fem}) sit squarely in the regime the criterion names:
multifractal states of effectively random couplings. Nature's free plate, whose genuine
true-operator effect is an order of magnitude weaker ($D_2 \approx 0.01$--$0.25$, ordered by
boundary adaptedness~\cite{ornelas2026fem}), correspondingly offers little of this resource ---
consistent with the campaign's central protocol-artifact finding. Enhancement exponents of
$1-D_2 \approx 0.24$--$0.44$ (AI band) versus $\approx 0.3$--$0.4$ (designed assemblies) are, on
this account, engineering targets rather than diagnostics.

Open, in order of leverage: (i) the Q2c size-tracking question above; (ii) energy-resolved
sharpening of the gasket null (registered as Q3c); (iii) the feasibility of an atomic
GOE$\to$GUE spacing measurement --- no atomic system has ever shown that crossover in level
statistics, and the mechanical version is now fully characterized
in-silico~\cite{ornelas2026fem} --- for which the kicked-rotor route remains the candidate
platform; and (iv) the standing scoop-watch on multifractal-sensing formulations, since the
ingredients of this composition are public and recent~\cite{sak2022,hu2026}.

\section{Methods and reproducibility}
\label{sec:methods}

Generalized-RP ensemble $H = \mathrm{diag}(\epsilon_i) + N^{-\gamma/2}W$, $\epsilon_i$ i.i.d.,
$W$ GOE; $\chif$ summed exactly over the spectrum per mid-window eigenstate; typical values are
log-averages; $D_2$ from IPR ladders over $N = 256$--$4096$ with window fractions held fixed.
Spacing statistics via the consecutive-spacing ratio~\cite{atas2013,oganesyan2007}, per symmetry
sector, against finite-size baselines computed with the identical protocol. Eigenpair
certification --- an independent no-vector eigenvalue reference, driver-level cross-checks,
agreement and residual gates per size --- is mandatory throughout; the one solver defect found
(\S\ref{sec:scope}) was caught by these gates and is documented with the fix. Every experiment
carries its preregistered reading in the script header; all numbers quoted above are frozen
verdicts. Compute: CPU-only commodity workstation.

\begin{thebibliography}{99}
\footnotesize
\bibitem{liu2014} W.-L.~Liu, X.~Xiong, L.~Song, X.~Wang (2014). Fidelity susceptibility and quantum
Fisher information for density operators with arbitrary ranks. \emph{Physica A} \textbf{410}, 167.
arXiv:1401.3154.
\bibitem{safranek2017} D.~\v{S}afr\'anek (2017). Discontinuities of the quantum Fisher information
and the Bures metric. \emph{Phys.\ Rev.\ A} \textbf{95}, 052320. arXiv:1612.04581.
\bibitem{hauke2016} P.~Hauke, M.~Heyl, L.~Tagliacozzo, P.~Zoller (2016). Measuring multipartite
entanglement through dynamic susceptibilities. \emph{Nature Physics} \textbf{12}, 778.
\bibitem{sak2022} M.~A.~Skvortsov, M.~Amini, V.~E.~Kravtsov (2022). Sensitivity of (multi)fractal
eigenstates to a perturbation of the Hamiltonian. \emph{Phys.\ Rev.\ B} \textbf{106}, 054208.
arXiv:2205.10297.
\bibitem{kravtsov2015} V.~E.~Kravtsov, I.~M.~Khaymovich, E.~Cuevas, M.~Amini (2015). A random
matrix model with localization and ergodic transitions. \emph{New J.\ Phys.} \textbf{17}, 122002.
\bibitem{mirlin1996} A.~D.~Mirlin, Y.~V.~Fyodorov, F.-M.~Dittes, J.~Quezada, T.~H.~Seligman
(1996). Transition from localized to extended eigenstates in the ensemble of power-law banded
random matrices. \emph{Phys.\ Rev.\ E} \textbf{54}, 3221.
\bibitem{sierpinski2025} Non-ergodic multifractal eigenstates from geometry alone: substitutional
defects on the Sierpi\'nski gasket. \emph{Commun.\ Phys.} (2025), s42005-025-02455-w.
arXiv:2410.18559.
\bibitem{tokarczyk2026} P.~Tokarczyk, L.~Vidmar, A.~Polkovnikov, P.~\L{}yd\.zba (2026). Fidelity
susceptibility across the Anderson transitions. arXiv:2605.25594.
\bibitem{hu2026} Hu et al.\ (2026). Preparation of critical/multifractal states in ultracold
atoms. arXiv:2605.21441.
\bibitem{sahoo2024} A.~Sahoo, U.~Mishra, D.~Rakshit (2024). Localization-driven quantum sensing.
\emph{Phys.\ Rev.\ A} \textbf{109}, L030601. arXiv:2305.02315.
\bibitem{sahoo2024b} A.~Sahoo, D.~Rakshit (2024). Enhanced sensing of Stark weak field under the
influence of Aubry--Andr\'e--Harper criticality. arXiv:2408.03232.
\bibitem{liu2024me} Y.~Liu et al.\ (2024). Fidelity-susceptibility metrology at exact mobility
edges in generalized AAH models. arXiv:2405.13282.
\bibitem{zapata2001} I.~Zapata, S.~Stringari et al.\ (2001). Wannier--Stark ladders and gravity.
\emph{Phys.\ Rev.\ A} \textbf{63}, 023607.
\bibitem{wilkinson1996} S.~R.~Wilkinson, C.~F.~Bharucha, K.~W.~Madison, Q.~Niu, M.~G.~Raizen
(1996). Observation of atomic Wannier--Stark ladders in an accelerating optical potential.
\emph{Phys.\ Rev.\ Lett.} \textbf{76}, 4512.
\bibitem{atas2013} Y.~Y.~Atas, E.~Bogomolny, O.~Giraud, G.~Roux (2013). Distribution of the ratio
of consecutive level spacings in random matrix ensembles. \emph{Phys.\ Rev.\ Lett.} \textbf{110},
084101.
\bibitem{oganesyan2007} V.~Oganesyan, D.~A.~Huse (2007). Localization of interacting fermions at
high temperature. \emph{Phys.\ Rev.\ B} \textbf{75}, 155111.
\bibitem{ornelas2026hyp} J.~A.~Ornelas Brand (2026). \emph{Why Rosenzweig--Porter? Evanescent
Coupling and a Random Matrix Hypothesis for Intermediate Statistics in Freely Vibrating Plates}.
Zenodo, v0.4.0. \href{https://doi.org/10.5281/zenodo.21313838}{doi:10.5281/zenodo.21313838}.
\bibitem{ornelas2026fem} J.~A.~Ornelas Brand (2026). \emph{Open-source finite-element tests of the
Rosenzweig--Porter hypothesis for freely vibrating plates: preliminary results}. Zenodo.
\href{https://doi.org/10.5281/zenodo.21314728}{doi:10.5281/zenodo.21314728}.
\bibitem{ornelas2026rmt} J.~A.~Ornelas Brand (2026). \emph{A random-matrix account of structure
formation in neural networks}. Zenodo, v0.1.0.
\href{https://doi.org/10.5281/zenodo.21314469}{doi:10.5281/zenodo.21314469}.
\end{thebibliography}

\end{document}
